\chapter{Layer 3: Refinement Validation with Liquid Haskell}
\label{ch:liquid}

\section{Purpose and Contract}
\label{sec:liq-purpose}

Layer~3 enforces \emph{structural} soundness before any proof obligation is
posed. Its inbound contract is ``a canonical IR expression''; its outbound
contract is either ``refinement-checked IR'' or a recorded refinement
violation. This is the first place where logic, not merely syntax, is checked.

\section{The Bridge}
\label{sec:liq-bridge}

The reference bridge (\texttt{layers/liquid/src/Bridge.hs}) imports the AST
and parser and exposes \texttt{generateLeanProof}:
\begin{lstlisting}[language=Haskell,caption={Liquid$\ensuremath{rightarrow}$Lean bridge (Bridge.hs).},label=lst:liq-bridge]
module MathLib5.Bridge where

import MathLib5.AST
import MathLib5.Parser

generateLeanProof :: Expr -> String
generateLeanProof expr =
  "theorem proof_obligation : " ++ toLean expr ++ " := by sorry"

toLean :: Expr -> String
toLean (IntVal i) = show i
toLean (Var v)    = v
toLean (App f xs) =
  "(" ++ toLean f ++ " " ++ unwords (map toLean xs) ++ ")"
toLean _          = "todo"
\end{lstlisting}

\section{What This Layer Actually Guarantees}
\label{sec:liq-guarantee}

In its present form the bridge \emph{generates} obligations rather than
discharging them: each expression becomes a Lean \texttt{theorem \ensuremath{dots} := by
sorry}, a placeholder to be closed at Layer~4 or Layer~10. The refinement
discipline is therefore \emph{staged}:
\begin{enumerate}
  \item L3 renders the expression into a Lean term and asserts it as a goal.
  \item L4 (Lean) attempts the proof; failure is recorded as a
        \texttt{sorry} in \texttt{sorrydb}.
  \item L10 (\texttt{sorryhunter}) later tries to close it.
\end{enumerate}
A full Liquid Haskell deployment would additionally carry SMT-decidable
refinement predicates (e.g., \texttt{\{e:Expr | wellFormed e\}}) so that
malformed expressions are rejected at L3 itself. The repository's
\texttt{metadata.json} for the layer records \texttt{status: "awaiting\_worm\_seal"},
indicating the refinement predicates are specified but not yet enforced in the
running pipeline.

\section{Refinement Predicates (Specified)}
\label{sec:liq-pred}

The intended refinement contract includes:
\begin{itemize}
  \item \texttt{wellFormed :: Expr -> Bool} --- no dangling variable
        references; all lambdas close over known bindings.
  \item \texttt{degreeNonNeg :: Poly -> Bool} --- polynomial degree $\ge 0$.
  \item \texttt{inBounds :: Index -> Vector -> Bool} --- indexing respects
        dimension.
\end{itemize}
These are the structural constraints the abstract promises but the current
code leaves as \texttt{todo} for non-literal cases. Honesty about this state
is part of the inspectability contract: a gap is a recorded gap, not a hidden
one.

\section{Connection to Layer 4}
\label{sec:liq-to-lean}

The crucial design point is that L3 is the \emph{generator} for L4. The
proof-obligations file produced here (\texttt{<base>.lean.obligations} in the
CI scaffold, \cref{ch:build}) is exactly the input Lean consumes. Thus the
boundary between refinement and proof is a file, not a function call --- a
clean, auditable seam.

\section{The Refinement Solved by the Prover}
\label{sec:liquid-solve}

For the summation fragment, the central Liquid predicate is
\[
\mathsf{sumOk}(v, r) \;\equiv\; r = \sum_{k=1}^{|v|} v_k.
\]
Liquid Haskell discharges \texttt{sumOk} at each call site using the
induction principle on the vector length, which the SMT solver accepts
because the arithmetic is linear. The proof obligation is generated
automatically; the programmer supplies only the type signature carrying
the refinement.
